% Generated draft: requires mathematical and source review.
\documentclass[11pt]{article}
\usepackage[margin=25mm]{geometry}
\usepackage{fontspec}
\setmainfont{Latin Modern Roman}
\usepackage{amsmath,amssymb,mathtools,unicode-math}
\setmathfont{Latin Modern Math}
\usepackage{enumitem,needspace,xurl,hyperref}
\setlength{\parindent}{0pt}
\setlength{\parskip}{5pt}
\setlength{\emergencystretch}{3em}
\newcommand{\N}{\mathbb{N}}
\newcommand{\Z}{\mathbb{Z}}
\newcommand{\Q}{\mathbb{Q}}
\newcommand{\R}{\mathbb{R}}
\newcommand{\C}{\mathbb{C}}
\newcommand{\F}{\mathbb{F}}
\newcommand{\A}{\mathbb{A}}
\newcommand{\PP}{\mathbb{P}}
\newcommand{\E}{\mathbb{E}}
\newcommand{\eps}{\varepsilon}
\newcommand{\dd}{\,\mathrm{d}}
\begin{document}
\section*{UnsolvedMath ID: 800030}
\textbf{Source problem code:} catalogue entry
\subsection{Definitions and mathematical statement}
A semigroup is a nonempty set \(S\) with an associative binary operation. It is finite if \(|S|<\infty\), and regular if
\[
\forall a\in S\;\exists b\in S\qquad aba=a.
\]
Write \(E(S)=\{e\in S:e^2=e\}\) for its set of idempotents. A regular semigroup is orthodox if \(E(S)\) is closed under multiplication. A band, or idempotent semigroup, is a semigroup satisfying \(x^2=x\). A variety of bands is a class of bands defined by semigroup identities.

\textbf{Source Problem 31 (S. I. Kublanovsky).} Consider the class
\[
\mathcal{O}_{\mathrm{fin}}
=\{S:S\text{ is a finite regular orthodox semigroup}\}.
\]
Is \(\mathcal{O}_{\mathrm{fin}}\) \emph{completely solvable}?

For a given variety \(\mathcal V\) of bands, consider also
\[
\mathcal{C}_{\mathcal V}
=\{S:S\text{ is a finite regular semigroup},\ E(S)\text{ is closed under multiplication},\ E(S)\in\mathcal V\},
\]
where \(E(S)\) carries the multiplication inherited from \(S\). Is \(\mathcal{C}_{\mathcal V}\) \emph{completely solvable}?

\textbf{Unresolved qualifications.} The supplied record does not define \emph{completely solvable} or specify which mathematical or decision problems this property concerns. Consequently, these questions preserve the source wording but do not yet specify a fully determined mathematical target. In the second question, interpreting the idempotents as belonging to a variety of idempotent semigroups requires that they form a subsemigroup under the inherited multiplication; this is the interpretation displayed above and requires source confirmation. The variety \(\mathcal V\) is treated as a given parameter. The record does not specify whether a uniform procedure over varieties, an effective presentation of \(\mathcal V\), or any restrictions on such varieties are intended.
\subsection{Short English statement}
Problem 31, attributed to S. I. Kublanovsky, asks whether finite regular orthodox semigroups are completely solvable, and asks the same question when their idempotent subsemigroups lie in a given variety of bands. The meaning of ``completely solvable'' is absent from the supplied record, so the intended target requires clarification from the source. The registry labels the problem open; that label is supplied metadata, not an independently verified assessment.
\subsection{Sources}
\begin{itemize}
\item Local UnsolvedMath registry, permanent ID \(800030\), source code \texttt{catalogue entry}, title \emph{The Lyapin's Notebook --- Problem 31}. The supplied record labels the problem \texttt{open} and was created and updated on 2026-10-16.
\item Bershadsky, Kublanovsky, and Mashevitzky, \emph{The Lyapin's Notebook: Unsolved Problems in Semigroup Theory} (2026), Problem 31. Supplied source URL: \url{https://arxiv.org/abs/2604.04763}. The record reports access on 2026-10-16 and extraction by \texttt{arxiv\_latex}; these are supplied provenance details, not a claim of independent source access or verification.
\item The problem statement attributes Problem 31 to S. I. Kublanovsky. The registry separately lists Bershadsky, Kublanovsky, and Mashevitzky in its proposer field, with proposed year 2026; that metadata does not independently establish the original proposal date or joint authorship of the problem.
\end{itemize}
\end{document}
